\usepackage[spanish,es-nodecimaldot]{babel}
\usepackage{amsmath,amssymb}
\usepackage{geometry}
\usepackage{microtype}
\usepackage{array}
\usepackage{booktabs}
\usepackage{xcolor}
\usepackage{hyperref}

\geometry{letterpaper,margin=2.5cm}

\hypersetup{
  colorlinks=true,
  linkcolor=blue!55!black,
  urlcolor=blue!55!black,
  citecolor=blue!55!black
}

\setlength{\parindent}{0pt}
\setlength{\parskip}{0.65em}
\setlength{\jot}{0.4em}

\newcommand{\curso}{F\'isica Aplicada a las Ciencias Farmac\'euticas}
\newcommand{\unidad}{Elementos de mec\'anica de la part\'icula en una dimensi\'on}
\newcommand{\vect}[1]{\boldsymbol{#1}}
\newcommand{\uvect}[1]{\hat{\boldsymbol{#1}}}
\newcommand{\R}{\mathbb{R}}
\newcommand{\dd}{\,\mathrm{d}}
\newcommand{\unidadSI}[1]{\,\mathrm{#1}}
\newcommand{\clase}[1]{\section*{#1}\addcontentsline{toc}{section}{#1}}
\newcommand{\subclase}[1]{\subsection*{#1}\addcontentsline{toc}{subsection}{#1}}

\newenvironment{objetivos}
  {\subclase{Objetivos de la clase}\begin{itemize}}
  {\end{itemize}}

\newenvironment{pizarra}
  {\subclase{Para desarrollar en pizarra}\begin{enumerate}}
  {\end{enumerate}}

\newenvironment{ejemplo}[1]
  {\subclase{Ejemplo: #1}}
  {}
